\chapter{Hỏi Khi Kiểm Tra Bài Cũ}
\chapterintro{Đừng để bài cũ chỉ thành chỗ bắt lỗi}{Một câu hỏi tốt có thể biến kiểm tra bài cũ thành lúc khởi động suy nghĩ.}

\section{Trong phần cũ, điều gì em nhớ mà không cần nhìn vở?}
\begin{cauhoibox}{Câu hỏi}
Trong phần cũ, điều gì em nhớ mà không cần nhìn vở?
\end{cauhoibox}
\begin{taisaocoich}
Câu hỏi này mở nhẹ hơn kiểu hỏi thẳng kiến thức chi tiết, nên lớp vào nhanh hơn và bớt căng.
\end{taisaocoich}
\begin{goiybien}
Sau khi nghe vài ý, mới đi vào câu hỏi cụ thể hơn nếu cần kiểm tra kỹ.
\end{goiybien}
\ngancach

\section{Nếu bạn quên bài, em sẽ nhắc bạn bằng câu nào?}
\begin{cauhoibox}{Câu hỏi}
Nếu bạn quên bài, em sẽ nhắc bạn bằng câu nào?
\end{cauhoibox}
\begin{taisaocoich}
Câu hỏi buộc học sinh diễn đạt lại kiến thức bằng lời của mình, rất tốt để đo mức hiểu thật.
\end{taisaocoich}
\begin{goiybien}
Có thể yêu cầu nói bằng đúng một câu, không dài quá 15 giây.
\end{goiybien}
\ngancach

\section{Phần nào của bài cũ đáng nhớ nhất chứ không chỉ đáng thuộc?}
\begin{cauhoibox}{Câu hỏi}
Phần nào của bài cũ đáng nhớ nhất chứ không chỉ đáng thuộc?
\end{cauhoibox}
\begin{taisaocoich}
Đây là câu hỏi kéo lớp sang mức hiểu sâu hơn, thay vì chỉ lặp lại nội dung máy móc.
\end{taisaocoich}
\begin{goiybien}
Hợp với các tiết ôn, tiết văn, giáo dục công dân, hoặc phần bài có ý nghĩa đời sống.
\end{goiybien}
\ngancach

\section{Nếu chỉ được nhắc bạn một từ về bài cũ, em chọn từ nào?}
\begin{cauhoibox}{Câu hỏi}
Nếu chỉ được nhắc bạn một từ về bài cũ, em chọn từ nào?
\end{cauhoibox}
\begin{taisaocoich}
Giới hạn vào một từ khiến học sinh phải cô ý rất nhanh. Câu hỏi này hợp để mở kiểm tra bài cũ theo cách ngắn mà có suy nghĩ.
\end{taisaocoich}
\begin{goiybien}
Sau các từ khóa, mời vài em giải thích ngắn vì sao mình chọn từ đó rồi mới đi sâu hơn.
\end{goiybien}
\ngancach

\section{Phần cũ nối được gì với bài hôm nay?}
\begin{cauhoibox}{Câu hỏi}
Phần cũ nối được gì với bài hôm nay?
\end{cauhoibox}
\begin{taisaocoich}
Câu hỏi này làm bài cũ bớt rời khỏi bài mới. Học sinh thấy việc nhắc lại không phải vì thủ tục, mà vì cần cho đoạn tiếp theo.
\end{taisaocoich}
\begin{goiybien}
Rất hợp khi muốn kiểm tra nhanh mà vẫn giữ được dòng chảy của tiết học.
\end{goiybien}
